\section*{الفصل \ref{ch:b3:partition-functions} --- الترموديناميك الإحصائي: دوال التجزئة}

\begin{solution}{exo:b3:partition-functions:1}
لنكتب كل توزيع على صورة الكمّات التي تحملها الهزازات الأربعة، من الأكبر
إلى الأصغر. «4، 0، 0، 0»: $4!/(3!\,1!) = 4$؛ «3، 1، 0، 0»: $4!/(2!\,1!\,1!) = 12$؛ «2، 2، 0، 0»:
$4!/(2!\,2!) = 6$؛ «2، 1، 1، 0»: $4!/(1!\,2!\,1!) = 12$؛ «1، 1، 1، 1»: 1. المجموع 35 $= \binom73$.
والأكثر احتمالًا هما «3، 1، 0، 0» و«2، 1، 1، 0» (12 لكل منهما)؛ والثاني، بإشغالاته
1 و2 و1 للكمّات 0 و1 و2، يتناقص منذ الآن مع الطاقة كما يفعل
\omterm{def:b3:partition-functions:boltzmann}{توزيع بولتزمان} من أجل الأعداد الكبيرة.
\end{solution}

\begin{solution}{exo:b3:partition-functions:2}
$\eu^{-2500/(8.314 \times 298.15)} = \eu^{-1.009} = 0.365$؛ وعند \qty{1000}{K}، $\eu^{-0.301} = 0.740$. وفي درجة حرارة
مرتفعة جدًا تؤول النسبة إلى 1 (إشغالان متساويان، ولا انقلاب في \omterm{def:b3:partition-functions:boltzmann}{الإشغال} أبدًا).
\end{solution}

\begin{solution}{exo:b3:partition-functions:3}
$m = 39.95\times10^{-3}/6.022\times10^{23} = \qty{6.634e-26}{kg}$؛ $\Lambda = h/\sqrt{2\pi mkT} =
\qty{1.600e-11}{m} = \qty{16.0}{pm}$؛ $q^{\mathrm T} = 10^{-3}/(1.600\times10^{-11})^3 = \num{2.44e29}$.
\end{solution}

\begin{solution}{exo:b3:partition-functions:4}
$hc/k = \qty{1.4388}{cm.K}$. \ce{N2}: $\theta_{\mathrm R} = \qty{2.863}{K}$، $q^{\mathrm R} = 298.15/(2 \times 2.863) = 52.1$.
\ce{HCl}: $\theta_{\mathrm R} = \qty{15.02}{K}$، $q^{\mathrm R} = 298.15/15.02 = 19.85$ (المجموع الدقيق 20.19).
\end{solution}

\begin{solution}{exo:b3:partition-functions:5}
\ce{I2}: $\theta_{\mathrm V} = 1.4388 \times 213.27 = \qty{306.9}{K}$، $q^{\mathrm V} = 1/(1 - \eu^{-1.029}) = 1.556$؛
والنسبة المثارة $1 - 1/q^{\mathrm V} = 0.357$. \ce{N2}: $\theta_{\mathrm V} = \qty{3352}{K}$، $q^{\mathrm V} = 1.000013$،
والنسبة المثارة \num{1.3e-5}.
\end{solution}

\begin{solution}{exo:b3:partition-functions:6}
\ce{NO}: $q^{\mathrm E} = 2 + 2\eu^{-119.73/207.22} = 2 + 2 \times 0.561 = 3.12$؛ ونسبة المستوى الأعلى $1.12/3.12
= 0.36$. وفي درجة حرارة مرتفعة $q^{\mathrm E} \to 4$ (مستويان متساويا \omterm{def:b3:partition-functions:boltzmann}{الإشغال}). \ce{Cl}:
$q^{\mathrm E} = 4 + 2\eu^{-882.35/207.22} = 4 + 2 \times 0.0142 = 4.03$.
\end{solution}

\begin{solution}{exo:b3:partition-functions:7}
$\ln q = \ln V + \frac32\ln T + \text{ثابت}$، ومنه $U - U(0) = NkT^2 \times \frac{3}{2T} = \frac32NkT$؛ ومن أجل
مول واحد $\frac32RT = \qty{3.72}{kJ/mol}$ عند \qty{298.15}{K}، و$C_V = \frac32R = \qty{12.47}{J.K^{-1}.mol^{-1}}$.
\end{solution}

\begin{solution}{exo:b3:partition-functions:8}
$kT/p^\circ = \qty{4.116e-26}{m^3}$، $\Lambda^3 = \qty{4.093e-33}{m^3}$، والنسبة \num{1.006e7}:
$S_{\mathrm m} = 8.3145 \times (16.124 + 2.5) = \qty{154.85}{J.K^{-1}.mol^{-1}}$، وهي القيمة الجدولية حتى
آخر رقم. وقسمة $p$ على 10 تضيف $R\ln10 = \qty{19.14}{J.K^{-1}.mol^{-1}}$.
\end{solution}

\begin{solution}{exo:b3:partition-functions:9}
$S^{\mathrm R}_{\mathrm m} = R[\ln(298.15/15.02) + 1] = 8.3145 \times (2.988 + 1) = \qty{33.16}{J.K^{-1}.mol^{-1}}$.
\end{solution}

\begin{solution}{exo:b3:partition-functions:10}
$\frac{\dd}{\dd J}\bigl[(2J+1)\eu^{-\theta J(J+1)/T}\bigr] = \bigl[2 - (2J+1)^2\theta/T\bigr]\eu^{-\theta J(J+1)/T} = 0$ تعطي
$2J + 1 = \sqrt{2T/\theta}$، أي $J_{\max} = \sqrt{T/2\theta} - \frac12$. \ce{HCl}: 2.66، ومنه $J = 3$؛ \ce{CO}:
6.86، ومنه $J = 7$. وشدة خط من \omterm{def:b3:rovibrational-spectroscopy:branches}{الفرع P} أو \omterm{def:b3:rovibrational-spectroscopy:branches}{الفرع R} تتبع
\omterm{def:b3:partition-functions:boltzmann}{إشغال} مستواه الأدنى، الذي تبلغ ذروته عند $J_{\max}$، على بعد بضعة خطوط من الفجوة
في مركز الحزمة.
\end{solution}

\begin{solution}{exo:b3:partition-functions:11}
$\int_0^\infty f\dd J = T/\theta$، $f(0) = 1$، $f'(0) = 2 - \theta/T$. ومنه $q^{\mathrm R} \approx T/\theta + \frac12 -
\frac16 + \frac{\theta}{12T} \approx T/\theta + \frac13$. والخطأ النسبي في $T/\theta$ نحو $\theta/3T$:
دون 1\,\% حين $T > 33\,\theta$. ومن أجل \ce{H2}، $\theta = 1.4388 \times 59.32 = \qty{85.3}{K}$: فعند
\qty{298.15}{K} يكون الخطأ 10\,\%، وهو غير مقبول (ويجب أيضًا أن يُعالَج قيد سبين النوى
على $J$ معالجة دقيقة، \cref{ch:b3:statistical-thermo-applied}).
\end{solution}

\begin{solution}{exo:b3:partition-functions:12}
من أجل $x = \theta_{\mathrm V}/T \to \infty$، $x/(\eu^x - 1) \to 0$ و$\ln(1 - \eu^{-x}) \to 0$: $S \to 0$. ومن أجل $x \to 0$،
$x/(\eu^x - 1) \approx 1 - x/2$ و$-\ln(1 - \eu^{-x}) \approx -\ln x + x/2$: $S \to R(1 - \ln x) = R[1 + \ln(T/\theta_{\mathrm V})]$.
\ce{I2}، $x = 1.029$: القيمة الدقيقة $R(0.5721 + 0.4420) = \qty{8.43}{J.K^{-1}.mol^{-1}}$؛ والنهاية $R(1 - 0.0289) =
\qty{8.07}{J.K^{-1}.mol^{-1}}$. فحتى عند $T \approx \theta_{\mathrm V}$ لا تقل النهاية الكلاسيكية إلا بنسبة 4\,\%.
\end{solution}

\begin{solution}{pb:b3:partition-functions:1}
\textbf{1.} $m = 28.014\times10^{-3}/6.02214\times10^{23} = \qty{4.652e-26}{kg}$.
\textbf{2.} $\Lambda = h/\sqrt{2\pi mkT} = \qty{1.910e-11}{m}$ (\qty{19.1}{pm}).
\textbf{3.} $kT/p^\circ = 1.380649\times10^{-23} \times 298.15/10^5 = \qty{4.116e-26}{m^3}$.
\textbf{4.} $q^{\mathrm T}/N = 4.116\times10^{-26}/(1.910\times10^{-11})^3 = \num{5.905e6}$: نحو ستة ملايين
حالة انسحابية لكل جزيء، فلا يكاد جزيئان يتقاسمان حالة واحدة، 
وتصح $Q = q^N/N!$.
\textbf{5.} $S^{\mathrm T}_{\mathrm m} = R(\ln\num{5.905e6} + \frac52) = 8.3145 \times 18.091 = \qty{150.42}{J.K^{-1}.mol^{-1}}$.
\textbf{6.} $+R\ln2 = \qty{5.76}{J.K^{-1}.mol^{-1}}$، أي \num{156.18}.
\textbf{7.} $B_0 = 1.998241 - 0.008659 = \qty{1.989582}{cm^{-1}}$.
\textbf{8.} $\theta_{\mathrm R} = 1.438777 \times 1.989582 = \qty{2.8626}{K}$.
\textbf{9.} $\sigma = 2$: النواتان متطابقتان؛ وقلب الجزيء طرفًا على طرف
يعطي اتجاهًا لا يمكن تمييزه، وتناظر التبادل للنواتين
لا يسمح، لكل حالة من حالات سبين النوى، إلا بنصف المستويات الدورانية.
\textbf{10.} $q^{\mathrm R} = 298.15/(2 \times 2.8626) = 52.08$.
\textbf{11.} نعم: $T/\theta_{\mathrm R} = 104$، أي خطأ قدره نحو 0.3\,\% في $q^{\mathrm R}$ وأقل بكثير في
الإنتروبيا.
\textbf{12.} $S^{\mathrm R}_{\mathrm m} = R(\ln52.08 + 1) = 8.3145 \times 4.9527 = \qty{41.18}{J.K^{-1}.mol^{-1}}$.
\textbf{13.} $J_{\max} = \sqrt{298.15/5.725} - 0.5 = 6.72$: $J = 7$ (إشغاله 0.0839، أعلى بقليل من
0.0831 للمستوى $J = 6$).
\textbf{14.} $2358.57 - 2 \times 14.324 = \qty{2329.92}{cm^{-1}}$.
\textbf{15.} $\theta_{\mathrm V} = 1.438777 \times 2329.92 = \qty{3352}{K}$.
\textbf{16.} $x = 3352.2/298.15 = 11.24$؛ $q^{\mathrm V} = 1/(1 - \eu^{-11.24}) = 1.000013$.
\textbf{17.} $\eu^{-11.24}/q^{\mathrm V} = \num{1.3e-5}$.
\textbf{18.} $S^{\mathrm V}_{\mathrm m} = R[11.24\eu^{-11.24} + \eu^{-11.24}] = \qty{0.0013}{J.K^{-1}.mol^{-1}}$، وهي مهملة.
\textbf{19.} $g_0 = 1$ (${}^1\Sigma$): $R\ln1 = 0$. والحالات الإلكترونية المثارة تعلو بعدة
إلكترون فولتات، أي أكثر من مئة مرة $kT$: فعوامل بولتزمان لها
مهملة تمامًا.
\textbf{20.} $150.42 + 41.18 + 0.00 + 0 = \qty{191.60}{J.K^{-1}.mol^{-1}}$.
\textbf{21.} الفرق \qty{0.01}{J.K^{-1}.mol^{-1}}، في مستوى
التقريبات المستعملة (\omterm{def:b3:quantum-model-systems:rotor}{الدوّار الصلب}، وفصل الحركات).
\textbf{22.} السعات الحرارية للصلب المقيسة انطلاقًا من بضعة كلفنات (والمستكملة
على صورة $T^3$ دونها)، وأنتالبيات التحول صلب--صلب والانصهار
والتبخر مقسومةً على درجات حرارتها، والسعة الحرارية للغاز،
وتصحيح للانحراف عن المثالية.
\textbf{23.} النوى محفوظة في كل تفاعل، فإنتروبيا سبينها هي
نفسها في الطرفين وتُلغى؛ والمسعرية لا تراها أيضًا، لأن
سبينات النوى تبقى غير منتظمة في البلورة حتى أدنى درجات الحرارة
المقيسة.
\textbf{24.} \textbf{$S^\circ(\ce{N2}, \qty{298.15}{K}) = \qty{191.6}{J.K^{-1}.mol^{-1}}$ من الثوابت الطيفية}،
منها 150.4 من الانسحاب و41.2 من الدوران.
\end{solution}
